% Complete Chapter 6 introduction, original PDF 73 / printed 63.
% block: exterior-calculus-chapter-heading page:73 type:chapter
\clearpage\setcounter{section}{5}
\section{{Basic notions of exterior calculus}}\label{chap:exterior-calculus}

% block: exterior-calculus-chapter-introduction page:73 type:ordinary-paragraph
{We have thus far studied Hilbert complexes in the abstract and mostly illustrated the concepts via the de Rham complex on a domain in three dimensions. In that example the complex has four spaces and three differential operators, which are the vector calculus operators gradient, curl, and divergence, whose properties we studied individually, employing very analogous arguments for each. In this second part of the book, we will focus on the de Rham complex in its most natural setting, on an $n$-dimensional manifold. The complex then has $n+1$ spaces connected by $n$ distinct operators. Using the language of exterior calculus we can introduce all these operators in a uniform way, obtaining a theory which is at once more elegant, more general, and more insightful. In this chapter, we present the relevant exterior calculus, culminating in the $L^2$ de Rham complex on a Riemannian manifold. Then, in the following chapter, we turn to the discretization of the corresponding Hodge Laplace boundary value problems using finite elements.}
